% Generated by scripts/materialize_unified_paper.py; do not edit.
\section{The exact-diagonal correlation certificate}\label{up:sec:exact}

We now prove the new analytic input.  Recall
\[
 u_0=\frac{1235783}{500000},\quad L=u_0^3,\quad a=e^{-u_0},
\]
and
\[
 E(z)=\sum_{n\ge0}\frac{z^n}{(3n)!},\quad
 H(z)=1+aE(z)^2,\quad G(z)=\frac{H(z)-H(-z)}2,
\]
\[
 F(x,y)=G(x)H(y)+G(y)H(x).
\]

\subsection{The bad-event separator}

For \(u\ge0\), define the cubic filters
\begin{align}
 P(u)&=\frac{e^u+2e^{-u/2}\cos(\sqrt3u/2)}3,\tag{U4.1}\label{up:main:eq:4-1}\\
 N(u)&=\frac{e^{-u}+2e^{u/2}\cos(\sqrt3u/2)}3.\tag{U4.2}\label{up:main:eq:4-2}
\end{align}
Then \(E(u^3)=P(u)\) and \(E(-u^3)=N(u)\).  Put
\[
 \sigma_0=10^{-8},\qquad T=2+2\sigma_0,qquad
 \sigma_*=\frac{1+2\sigma_0}{2}=\frac{50000001}{100000000}.
 \tag{U4.3}\label{up:main:eq:4-3}
\]

\begin{lemma}[Half-line ratio]\label{up:lem:ratio}
For every \(u\ge u_0\),
\[
 \frac{1+aP(u)^2}{1+aN(u)^2}>T.
 \tag{U4.4}\label{up:main:eq:4-4}
\]
\end{lemma}

\begin{proof}
Let \(R(u)=(1+aP(u)^2)/(1+aN(u)^2)\).  Its derivative has the sign of
\[
 W(u)=PP'(1+aN^2)-NN'(1+aP^2).
 \tag{U4.5}\label{up:main:eq:4-5}
\]
An outward-rounded enclosure proves
\(R(u_0)-T>2.6595\cdot10^{-7}\) and \(W>0\) on 65,536 rational cells
covering \([u_0,2.9]\).

For \(u\ge2.9\), use
\[
 P(u)\ge\frac{e^u-2e^{-u/2}}3,qquad
 |N(u)|\le\frac{2e^{u/2}+e^{-u}}3.
\]
After clearing the positive denominator, it suffices to prove positivity of
\[
 B_T(u)=e^{2u}-4Te^u-4e^{u/2}-4Te^{-u/2}+4e^{-u}-Te^{-2u}.
\]
At \(u=2.9\), interval arithmetic proves
\(aB_T-9(T-1)>0\).  Moreover
\[
 B_T'(u)\ge D(u):=2e^u(e^u-2T)-2e^{u/2}-4e^{-u},
\]
and
\[
 D'(u)=4e^{2u}-4Te^u-e^{u/2}+4e^{-u}>0
\]
throughout the half-line by exponential dominance, with a strict endpoint
check.  Hence \(B_T\), and therefore \(R-T\), remain positive.
\end{proof}

\begin{lemma}[Separator]\label{up:lem:separator}
If one of \(A_1,A_2\) is less than \(-1\), then
\[
 F(LA_1,LA_2)<-\sigma_*.
 \tag{U4.6}\label{up:main:eq:4-6}
\]
\end{lemma}

\begin{proof}
Let \(x=-t\) with \(t\ge L\), put \(A=H(-t)\), and set
\(R=H(t)/H(-t)\).  The exact identity
\[
 F(-t,y)=\frac A2\bigl((2-R)H(y)-H(-y)\bigr)
 \tag{U4.7}\label{up:main:eq:4-7}
\]
holds for every real \(y\).  By \cref{up:lem:ratio}, \(R>2+2\sigma_0\),
whereas \(A,H(y),H(-y)\ge1\).  Thus the right side of \eqref{up:main:eq:4-7} is less than
\(-(1+2\sigma_0)/2=-\sigma_*\).  Exchange the coordinates if necessary.
\end{proof}

\subsection{Complete two-dimensional reduction}

\begin{lemma}[Diagonal reduction]\label{up:lem:diagonal-reduction}
If \(x,y\ge-L\) and \(m=\max\{x,y,0\}\), then
\[
 F(x,y)\le F(m,m).
 \tag{U4.8}\label{up:main:eq:4-8}
\]
\end{lemma}

\begin{proof}
The series for \(H\) has nonnegative coefficients, so \(H\), and therefore
its odd part \(G\), are nondecreasing on the positive half-line.  If
\(x=-v^3\in[-L,0]\), the root-filter formula gives
\[
 |E(-v^3)|\le e^{v/2},\qquad
 H(x)\le1+e^{-u_0}e^v\le2.
 \tag{U4.9}\label{up:main:eq:4-9}
\]
Assume \(x\le y\).  If \(y<0\), both odd parts are negative and hence
\(F(x,y)<0=F(0,0)\).  If \(x<0\le y\), discard the nonpositive summand and
use \eqref{up:main:eq:4-9}:
\[
 F(x,y)\le G(y)H(x)\le2G(y)\le2G(y)H(y)=F(y,y).
\]
If \(0\le x\le y\), monotonicity bounds both summands by \(G(y)H(y)\).
This exhausts all sign patterns and their boundaries.
\end{proof}

On the diagonal we retain the exact identity
\[
 F(z,z)=2G(z)H(z)=H(z)^2-H(z)H(-z),\qquad z\ge0.
 \tag{U4.10}\label{up:main:eq:4-10}
\]

\subsection{Compact interval and analytic tail}

Let \(z=u^3\) and \(w=(u^3+u_0^3)^{1/3}\).  Define
\[
 D_0=\frac{88053}{10^9}.
\]

\begin{lemma}[Exact diagonal envelope]\label{up:lem:envelope}
For every \(u\ge0\),
\[
 \frac{H(u^3)^2-H(u^3)H(-u^3)}{e^{4w}}<D_0.
 \tag{U4.11}\label{up:main:eq:4-11}
\]
\end{lemma}

\begin{proof}
The compact interval \([0,20]\) is covered by 131,072 exact rational cells.
On each cell a 512-bit Arb evaluation encloses the left side of \eqref{up:main:eq:4-11}; the
minimum slack is greater than \(4.2687\cdot10^{-6}\).

For \(u\ge20\), use \(H(-u^3)\ge1\) and write
\(H(u^3)=1+aP(u)^2\).  Then
\[
 H(u^3)^2-H(u^3)H(-u^3)
 \le aP(u)^2+a^2P(u)^4.
\]
Since
\[
 |P(u)|\le\frac{e^u}{3}(1+2e^{-3u/2})
\]
and \(w\ge u\), the normalized expression is at most
\[
 \frac a9e^{-2u}(1+2e^{-3u/2})^2
 +\frac{a^2}{81}(1+2e^{-3u/2})^4.
 \tag{U4.12}\label{up:main:eq:4-12}
\]
Both terms decrease.  At \(u=20\), \eqref{up:main:eq:4-12} is below
\(8.805216326983\cdot10^{-5}\), leaving strict slack greater than
\(8.3673\cdot10^{-10}\); its limiting value \(a^2/81\) has the same strict
ordering.  This proves the entire half-line.
\end{proof}

Combining \cref{up:lem:diagonal-reduction,up:lem:envelope}, if
\(A_1,A_2\ge-1\) and \(M=\max(A_1,A_2)\), then
\[
 F(LA_1,LA_2)<D_0\exp\!\left(4u_0(M+1)^{1/3}\right).
 \tag{U4.13}\label{up:main:eq:4-13}
\]

\subsection{Exact moment/tail budget}

The final passage from the separator to a correlation tail is short but
proof-critical, so we give it in full rather than treating it as an unnamed
interface.

\begin{lemma}[Positive-moment expectation]\label{up:lem:positive-expectation}
For any finite set \(\mathcal X\), independent identically distributed
\(U,U'\in\mathcal X\), real Hilbert spaces \(\mathcal H_1,\mathcal H_2\),
and maps \(\sigma_i:\mathcal X\to\mathcal H_i\), put
\(Z_i=\langle\sigma_i(U),\sigma_i(U')\rangle\).  Then
\[
 \mathbb E F(LZ_1,LZ_2)\ge0.
 \tag{U4.14}\label{up:eq:positive-expectation}
\]
\end{lemma}

\begin{proof}
The Taylor coefficients of \(E\) and \(H\) are nonnegative, and taking the
odd part of \(H\) merely retains its odd-degree coefficients.  Hence every
bivariate Taylor coefficient of
\(F(x,y)=G(x)H(y)+G(y)H(x)\) is nonnegative.  For nonnegative integers
\(r,s\), tensorization gives
\[
 \mathbb E(Z_1^rZ_2^s)
 =\left\|\mathbb E\bigl(
   \sigma_1(U)^{\otimes r}\otimes
   \sigma_2(U)^{\otimes s}\bigr)\right\|^2\ge0.
\]
Because \(\mathcal X\) is finite, \((Z_1,Z_2)\) has bounded support; the
entire Taylor series is therefore absolutely integrable and may be averaged
term by term.  Summing the nonnegative terms proves
\eqref{up:eq:positive-expectation}.
\end{proof}

Take
\[
 \beta=\frac{330867}{500000},\qquad
 \varepsilon=\frac{6893}{10^7}.
\]
Exact rational arithmetic gives
\[
 \sigma_*(1-\beta)-D_0\beta(1+2/\varepsilon)
 =\frac{39437634699447}{3446500000000000000}>10^{-5}.
 \tag{U4.15}\label{up:eq:tail-budget}
\]

\begin{proposition}[Certified correlation tail]\label{up:prop:correlation-tail}
The exact separator, envelope, and budget above imply
\[
 \cG_2^{(3)}\!\left(
 \frac{330867}{500000},
 4u_0(1+\varepsilon)\right).
\]
\end{proposition}

\begin{proof}
Fix data from the definition \eqref{up:eq:correlation-definition}, put
\(\mathcal E=\{Z_1\ge-1,Z_2\ge-1\}\),
\(M=\max\{Z_1,Z_2\}\), \(Y=((M+1)_+)^{1/3}\), and
\(c=4u_0\).  By \cref{up:lem:positive-expectation,up:lem:separator} and
\eqref{up:eq:positive-expectation}, while \eqref{up:main:eq:4-13} controls the good event,
\[
 0\le\mathbb EF(LZ_1,LZ_2)
 <D_0\mathbb E(e^{cY}\mathbf1_{\mathcal E})
   -\sigma_*\mathbb P(\mathcal E^{\mathrm c}).
 \tag{U4.16}\label{up:eq:master-tail}
\]
If \(\mathbb P(\mathcal E)\ge\beta\), choose \(\lambda=-1\) and either
coordinate in \eqref{up:eq:correlation-definition}.  Otherwise suppose, for a
contradiction, that the claimed correlation conclusion fails for both
coordinates at every \(\lambda\ge-1\).  For \(v\ge0\), taking
\(\lambda=v^3-1\) and using the union bound gives
\[
 \mathbb P(\mathcal E\cap\{Y\ge v\})
 <2\beta e^{-c(1+\varepsilon)v}.
 \tag{U4.17}\label{up:eq:Y-tail}
\]
The layer-cake identity and Tonelli's theorem now give
\begin{align*}
 \mathbb E(e^{cY}\mathbf1_{\mathcal E})
 &=\mathbb P(\mathcal E)
   +c\int_0^\infty e^{cv}
       \mathbb P(\mathcal E\cap\{Y>v\})\,dv\\
 &<\beta+2\beta c\int_0^\infty e^{-c\varepsilon v}\,dv
 =\beta(1+2/\varepsilon).
\end{align*}
Substitution into \eqref{up:eq:master-tail} makes its right-hand side at most
\(D_0\beta(1+2/\varepsilon)-\sigma_*(1-\beta)<0\), by
\eqref{up:eq:tail-budget}.  This contradiction proves the proposition.
\end{proof}

Finally,
\[
 4u_0(1+\varepsilon)
 =\frac{12366348252219}{1250000000000}
 =9.8930786017752,
\]
which is the exact handoff asserted in \eqref{up:eq:correlation-input}.
